\originalchapter{Lyman吸收森林与星系际气体}
本章对应8.902第21讲。远方类星体提供亮的背景光，沿途气体在自己的静止系吸收特定波长；我们因此能从一条观测光谱识别不同红移处的气体。吸收线既依赖密度，又依赖电离、温度和速度。本章从频率相关截面推导光学深度，并说明Gunn--Peterson暗槽能告诉我们什么以及不能单独告诉什么。

\originalsection{为什么森林位于类星体Ly\texorpdfstring{$\alpha$}{alpha}发射线的蓝侧}
氢$1s\leftrightarrow2p$的Lyman-$\alpha$静止波长约$\lambda_\alpha=1215.67$埃。红移$z_q$的类星体发射线位于$\lambda_\alpha(1+z_q)$；沿途红移$z_a<z_q$的吸收体在$\lambda_{ \rm obs}=\lambda_\alpha(1+z_a)$处吸收，因此该线在类星体发射线的蓝侧。光子发射时可来自类星体更短波长的\emph{连续谱}，并非必须先在类星体处是一个Ly$\alpha$线光子，再沿途被同一条线重复吸收。

一连串不同$z_a$的气体区域产生森林。谱中的金属双线还可帮助识别独立吸收系统，例如Mg II的2796、2803埃双线，但一条森林线不能仅凭深浅唯一指定一个孤立球云。实际星系际气体沿宇宙网连续分布，本动速度、温度与密度共同塑造观测线形。

\originalsection{柱密度与三个吸收尺度}
中性氢柱密度$N_{\rm HI}=\int n_{\rm HI}\dd l$的单位为$\mathrm{cm}^{-2}$。在Lyman极限912埃附近，氢光电离截面约$6.3\times10^{-18}\ \mathrm{cm}^2$，故$N_{\rm HI}\sim1.6\times10^{17}\ \mathrm{cm}^{-2}$已有连续吸收光学深度约1。较低柱密度通常表现为森林线，足够高时成为Lyman极限系统。

常用阻尼Ly$\alpha$系统定义$N_{\rm HI}\ge2\times10^{20}\ \mathrm{cm}^{-2}$，高柱密度使自然线宽的翼显著。原课印成$10^{32}\ \mathrm{cm}^{-2}$，不能原样用作物理分类。该定义另核对\href{https://arxiv.org/abs/astro-ph/0509481}{Wolfe、Gawiser与Prochaska（2005）}的开篇分类；只引用定义，不重印受限综述或其图。柱密度分类是观测约定，具体线宽仍受气体速度结构影响。

\originalsection{从归一线型到光学深度}
本节用Gaussian电磁单位。某共振线的吸收截面为
\[
 \sigma_\nu=S_\alpha\varphi(\nu-\nu_\alpha),\qquad
 S_\alpha=\frac{\pi e^2}{m_ec}f_\alpha,\qquad\int\varphi(\nu)\dd\nu=1.
\]
振子强度$f_\alpha$是无量纲矩阵元参数，不是逐次吸收概率。$\varphi$单位为$\mathrm{Hz}^{-1}$，$S_\alpha$单位为$\mathrm{cm}^2\mathrm{Hz}$，乘积$\sigma_\nu$才有面积单位。Ly$\alpha$的$f_\alpha\simeq0.416$给$S_\alpha\simeq1.10\times10^{-2}\ \mathrm{cm}^2\mathrm{Hz}$；原课将系数与线型的单位混在一起，应改正。

热速度使线心近Gaussian，自然有限寿命给Lorentz翼，两者卷积为Voigt型。热Doppler参数$b=\sqrt{2\kb T/m_H}$，$\Delta\nu_D=(b/c)\nu_\alpha$；若另有Gaussian非热速度，平方宽度相加。因此$b=20\ \mathrm{km\,s^{-1}}$在\emph{纯热}假设下给$T\simeq2.4\times10^4$ K，非热展宽存在时给的是温度上限。Lorentz翼按频差的平方反比衰减，不能把整个Voigt型都称为Gaussian。

沿观测频率$\nu_o$的光路，当地光子频率是$\nu_o(1+z)$，当地物理路径元为
\[
 \dd l=c|\dd t|=\frac{c\dd z}{(1+z)H(z)}.
\]
因此均匀背景近似下
\begin{equation}\label{eq:astro-lya-opacity}
 \tau(\nu_o)=\int_0^{z_q}\frac{c\,n_{\rm HI}(z)}{(1+z)H(z)}
 S_\alpha\varphi[\nu_o(1+z)-\nu_\alpha]\dd z.
\end{equation}
$n_{\rm HI}$是\emph{物理}数密度。射出原窄光束的共振散射也表现为消光，若忽略重新散入该光束的光，$F_\nu=F_{\nu,\rm cont}\ee^{-\tau}$。这不是将未经距离和红移变换的源静止系$F$直接代到观测谱中。

\originalsection{Gunn--Peterson光学深度与饱和的限制}
若线型很窄、沿共振区域的背景变化缓慢，用$\varphi\simeq\delta_D$，令$z_r=\nu_\alpha/\nu_o-1$。Dirac变量变换给$\delta_D[\nu_o(1+z)-\nu_\alpha]=\delta_D(z-z_r)/\nu_o$，代入式\eqref{eq:astro-lya-opacity}，有
\begin{equation}\label{eq:astro-gp-opacity}
 \tau_{ \rm GP}(z_r)=\frac{S_\alpha\lambda_\alpha n_{\rm HI}(z_r)}{H(z_r)}.
\end{equation}
在物质主导期$H\simeq H_0\sqrt{\Omega_{m,0}}(1+z)^{3/2}$，若中性比例$x_{\rm HI}$不变而总氢密度随$(1+z)^3$增长，便有$\tau_{ \rm GP}\propto x_{\rm HI}(1+z)^{3/2}$。注意$S_\alpha\lambda_\alpha$单位是$\mathrm{cm}^3\mathrm{s}^{-1}$，所以最后结果无量纲。

课程用$z=3$、$n_H\sim10^{-5}\ \mathrm{cm}^{-3}$、$H_0\sim70\ \mathrm{km\,s^{-1}\,Mpc^{-1}}$，完全中性时得到$\tau\sim10^5$。要透过显著光，气体的中性比例可需低到$10^{-5}$量级。反过来，$x_{\rm HI}$即使远小于1也能造成饱和黑暗槽。\textbf{黑暗槽不能单独证明该处气体“大部分中性”}，它往往只给中性比例的下限；课程以$z\sim6$暗槽直接推断之前大部分中性，需保留这项限制。定量再电离分析还用阻尼翼、空间起伏和其他数据。

\originalsection{光电离如何维持少量中性氢}
若背景紫外辐射场的角平均强度为$J_\nu$，每个中性原子的光电离率为
\[
 \Gamma_{ \rm HI}=4\pi\int_{\nu_{ \rm LL}}^\infty
 \frac{J_\nu\sigma_{ \rm ion}(\nu)}{h\nu}\dd\nu.
\]
它单位为$\mathrm{s}^{-1}$，不是整个宇宙每秒仅发出多少光子的总率。忽略碰撞电离、取适当复合系数$\alpha(T)$时，局部电离平衡要求
\[
 \Gamma_{ \rm HI}n_{\rm HI}=\alpha(T)n_en_{\rm HII}.
\]
高度电离氢气$n_e\simeq n_{\rm HII}\simeq n_H$，于是$x_{\rm HI}\simeq\alpha n_H/\Gamma_{ \rm HI}$。以课程例$\Gamma\sim10^{-12}\ \mathrm{s}^{-1}$、$\alpha\sim4\times10^{-13}\ \mathrm{cm^3\,s^{-1}}$、$n_H\sim10^{-5}\ \mathrm{cm}^{-3}$，给$x_{\rm HI}\sim4\times10^{-6}$，与透明度要求相容。

每个中性原子的电离时间约$\Gamma^{-1}\sim3\times10^4$年，而每个离子的复合时间$1/(\alpha n_e)\sim2.5\times10^{17}$秒，约$8\times10^9$年。原课将$2\times10^{17}$秒换成$3\times10^6$年是单位换算错误。二者差异由平衡时\emph{中性原子很少}调节，不能靠把两个不同比例的粒子时间强行设成相同而推导平衡。

线宽对应的$\kb T$小于13.6 eV，不足以单独排除所有碰撞电离，因为Maxwell分布仍有高能尾部；必须比较碰撞电离与复合率。所需极低中性比例和强紫外背景使光电离成为低密度森林气体的合适解释。复合系数用Case A还是Case B取决于再发出的电离光子是否逃逸；气体团聚又使$\langle n_e n_p\rangle$不同于平均密度平方。

\originalsection{光谱能推断什么，以及还缺什么}
由$x_{\rm HI}\propto\alpha(T)n_H/\Gamma$和式\eqref{eq:astro-gp-opacity}可看出，吸收深度对密度近乎二次敏感，也依赖紫外背景和温度。相同$\tau$可以来自不同密度与电离条件；线宽有热和速度简并，饱和线又失去线心柱密度信息。金属双线、氘线、阻尼翼和不同Lyman线可帮助打破部分简并。

从森林通量谱推到物质功率谱，还需非线性流体、热史、速度畸变和辐射转移模型。原课提到的THESAN和HERA说明研究方法，但本卷没有运行相应模拟或处理实际望远镜数据，不能把课堂示意图当作已完成的模型核验。星系际介质、再电离及暗物质结构因此通过同一条吸收光路联系起来，又保留明确的反演边界。
